% =============================================================
%  ch25_hebelarm.tex  --  Das Hebelarm- und Schwingungsproblem
% =============================================================

\chapter{Das Hebelarm- und Schwingungsproblem}

Der Arm hängt mit Abstand $L$ unter der X-Achse. Bei harten Beschleunigungen
zur Zielposition entstehen Reaktionskräfte und -momente, die den Ausleger zum
Schwingen anregen.

\section{Die Physik}
Beschleunigt der Schlitten mit $a$, wirkt auf die bewegte Masse $m$ die
Trägheitskraft $F=m\,a$. Über den Hebelarm $L$ ergibt das am Anbindungspunkt
ein \textbf{Biegemoment}
\begin{equation}
M = F\cdot L = m\,a\,L.
\end{equation}
Der Ausleger verhält sich näherungsweise wie ein einseitig eingespannter Balken
(Kragträger) mit Steifigkeit $k=\dfrac{3EI}{L^3}$ und einer ersten
\textbf{Eigenfrequenz}
\begin{equation}
f_1 = \frac{1}{2\pi}\sqrt{\frac{k}{m_{\text{eff}}}}
    = \frac{1}{2\pi}\sqrt{\frac{3EI}{m_{\text{eff}}\,L^3}}.
\end{equation}
$E$ = E-Modul, $I$ = Flächenträgheitsmoment des Auslegerquerschnitts.
Wichtig: $f_1$ fällt mit $L^{3/2}$ -- ein längerer Ausleger ist drastisch
weicher. Anregungsanteile der Bewegung nahe $f_1$ führen zu Resonanz und langem
Nachschwingen, was die Schlaggenauigkeit ruiniert.

\section{Gegenmaßnahmen}
\begin{enumerate}
\item \textbf{Bewegte Masse minimieren:} Motoren rahmenfest (Riemen-/Seilzug
  zu den Armgelenken, CoreXY-Logik), Arm aus Aluminium/CFK, Schwerpunkt nah an
  der Achse.
\item \textbf{Steifigkeit maximieren:} große Stab-/Schienenquerschnitte
  (gestützte Linearschienen statt dünner Stäbe), Doppelführung, kurzer
  Ausleger, triangulierter Rahmen. $f_1$ deutlich \emph{über} die
  Betriebsfrequenzen legen.
\item \textbf{Ruckbegrenzte Bewegungsprofile (S-Kurve):} sanftes Beschleunigen
  regt hohe Frequenzen weniger an.
\item \textbf{Input Shaping / Resonanzkompensation:} der Steuerbefehl wird so
  vorgeformt, dass die Anregung bei $f_1$ ausgelöscht wird -- genau das, was
  Bambu/Klipper mit einem Beschleunigungssensor (ADXL345) als ``resonance
  compensation'' tun. Das ist die wirksamste softwareseitige Maßnahme.
\item \textbf{Dämpfung:} Strukturdämpfung, ggf.\ ein abgestimmter Tilger
  (tuned mass damper); Riemenvorspannung erhöhen (Riemendehnung ist Teil der
  Nachgiebigkeit).
\end{enumerate}

\begin{infobox}[title=Anknüpfung an den FEM-Hintergrund]
Bestimme $f_1$ und die Modenformen per \textbf{Modalanalyse} (SolidWorks
Simulation / FEM). Das liefert dir genau die Frequenz, auf die du das Input
Shaping abstimmst, und zeigt, welche Strukturteile du versteifen musst. Das
verbindet FEM-Erfahrung direkt mit der Reglerauslegung.
\end{infobox}
